% =========================================================
% setup/00_metadata.tex
% User-editable report metadata
% =========================================================
%
% PURPOSE
% -------
% This is the primary configuration file for a new report.
% In most projects, users should edit this file first and leave
% main.tex and frontmatter/titlepage.tex unchanged.
%
% TEXT SAFETY
% -----------
% Escape LaTeX special characters in metadata when needed:
%   &  -> \&        %  -> \%        _  -> \_
% Use plain text in fields that are copied into PDF metadata.
%

% =========================================================
% 1. Course, program, or project context
% =========================================================

\newcommand{\coursenumber}{ENGR 6000}
\newcommand{\coursetitle}{Advanced Engineering Analysis}
\newcommand{\coursename}{\coursenumber\ \coursetitle}
\newcommand{\courseshorttitle}{\coursenumber\ Engineering Analysis}

% Use these fields for an academic department and institution,
% a company and business unit, or a research group and laboratory.
\newcommand{\departmentname}{Department of Electrical and Computer Engineering}
\newcommand{\universityname}{Example University}

% =========================================================
% 2. Report identity
% =========================================================

\newcommand{\assignmenttype}{Technical Project Report}
\newcommand{\projecttitle}{Smart Sensor Array Design and Validation}
\newcommand{\projectshorttitle}{Smart Sensor Array}

% Formatted title shown on the cover page. Line breaks are intentional.
\newcommand{\reporttitle}{%
    \assignmenttype:\\[0.35cm]
    Smart Sensor Array Design\\[0.15cm]
    and Validation%
}

% Plain-text title used for PDF properties and accessibility metadata.
% Keep this version free of formatting commands and manual line breaks.
\newcommand{\reporttitleplain}{\assignmenttype: \projecttitle}

% Short report title shown in page headers.
\newcommand{\reportshorttitle}{\projectshorttitle}

% =========================================================
% 3. Author, supervisor, and submission information
% =========================================================

\newcommand{\studentname}{Author Name}
\newcommand{\studentid}{000000000}
\newcommand{\studentidlabel}{Student ID}
\newcommand{\professorname}{Dr. Faculty Name}
\newcommand{\semestername}{Fall Semester 2026}
\newcommand{\submissiondate}{\today}

% Optional title-page rows. Set each switch to true or false.
\newif\ifshowstudentid
\showstudentidtrue

\newif\ifshowsupervisor
\showsupervisortrue

\newif\ifshowsemester
\showsemestertrue

% Labels may be adapted for academic, industrial, or consulting reports.
\newcommand{\authorlabel}{Author}
\newcommand{\supervisorlabel}{Instructor / Supervisor}

% =========================================================
% 4. Running headers
% =========================================================

\newcommand{\headerleft}{\courseshorttitle}
\newcommand{\headerright}{\reportshorttitle}

% =========================================================
% 5. Optional cover logo
% =========================================================

% Set \showreportlogotrue to display the logo and
% \showreportlogofalse to suppress it without deleting the path.
\newif\ifshowreportlogo
\showreportlogotrue

% Place the authorized image inside report-template/figures/.
\newcommand{\reportlogo}{figures/logo-placeholder.png}
\newcommand{\reportlogowidth}{7.2cm}

% The template intentionally renders a visible warning box when the
% logo switch is enabled but the configured image cannot be found.

% =========================================================
% 6. PDF document properties
% =========================================================

\newcommand{\pdfsubjecttext}{\coursename}
\newcommand{\pdfkeywordstext}{engineering report, modeling, simulation, validation, signal processing, embedded systems}
